
\needspace{4\baselineskip}
\color{uniblue}{\subsection{Блокировка при качаниях (БК)}}
\color{black}

\begin{enumerate}[label=\arabic{section}.\arabic{subsection}.\arabic*, labelsep=4pt, leftmargin=0pt, itemindent=57pt]

\item
Блокировка при качаниях используется совместно с функцией дистанционной защиты для предотвращения ложного срабатывания ступеней ДЗ в режиме качаний мощности и при асинхронном ходе. Работа ступеней ДЗ вводится при КЗ на время с последующим выводом. Измерительные органы блокировки при качаниях реагируют на приращение во времени векторов тока прямой и обратной последовательности.

\item
В устройстве блокировка при качаниях представлена функцией «БК».

\item
Функция имеет два органа:
\begin{itemize}
\item чувствительный;
\item грубый.
\end{itemize}

\item
Основные характеристики функции:
\begin{itemize}
    \item основная погрешность срабатывания ИО тока по приращению не более 15\%;
    \item время срабатывания ИО по приращению тока прямой и обратной последовательности при условии повышения тока со скоростью равной $2dI_{\textsl{уст.}}$ не более 25 мс.
\end{itemize}

\item
Алгоритм функции по скорости изменения тока, включает в себя:
\begin{itemize} 
    \item двухступенчатый ИО, реагирующий на скачкообразное изменение тока обратной последовательности; 
    \item двухступенчатый ИО, реагирующий на скачкообразное изменение тока прямой последовательности. 
\end{itemize}

\item
Принцип работы блокировки при качаниях  основан на двух основных признаках режима качаний, позволяющих отличить его от режима КЗ:
\begin{itemize}  
    \item в режиме качаний изменение фазных токов носит плавный характер; 
    \item качания является симметричным режимом  (даже при трехфазном КЗ имеет место кратковременная несимметрия).
\end{itemize}

\item
Измерительными  органами  функции  являются  реле,  реагирующие  на  приращения  тока обратной последовательности, обеспечивающие работу функции при всех видах несимметричных КЗ. Появление  кратковременной  несимметрии  перед  трехфазным  КЗ  достаточно,  чтобы  обеспечить блокировку  и  при  симметричном  КЗ.  Несмотря  на  это,  в  алгоритм  работы  функции дополнительно введены ИО, реагирующие на приращения тока прямой последовательности.

\item
В  функции  формируются  два  выходных  сигнала:  «БК:  Срабатывание  БК-б»  (ввод быстродействующих ступеней защит) и «БК: Срабатывание БК-м» (ввод медленнодействующих ступеней). Под быстродействующими понимают ступени ДЗ, имеющие выдержку времени на срабатывание меньше периода возможных в системе качаний, вызванных внешним КЗ. Обычно период качаний составляет 1,5–2,0 с.  Ступени,  имеющие  выдержку  времени  на  срабатывание  больше  периода  качаний,  называют медленнодействующими. 

\item
Функция  содержит  в  себе  два  ИО  приращения  тока  прямой  последовательности  (грубый  и чувствительный) и два ИО приращения тока обратной последовательности (грубый и чувствительный). С помощью уставки программного переключателя SGF2, в функции предусмотрен ускоренный возврат логики БК по наличию сигнала «В Отключен» на случай включения на КЗ в режиме АПВ. 

\item
В  неаварийном  режиме  качаний  ИО  сопротивления  могут  сработать  ложно.  При  этом  ИО  БК, отстроенные  от  режима  качаний  выбором  уставок  по  изменению  токов  прямой  и  обратной последовательностей, не срабатывают и не пускают ДЗ. В случае возникновения КЗ срабатывают ИО БК, которые  разрешают  прохождение  сигналов  срабатывания  от  ИО  ступеней  на  время,  определяемое уставкой  Т1,  при  срабатывании  чувствительного  канала  или  на  время,  определяемое  уставкой  Т2,  при срабатывании грубого канала. Полный возврат функции БК происходит по истечение выдержки времени Т3.  В  случае  пуска  алгоритма  от  грубых  ступеней  ИО  возможен  повторный  ввод быстродействующих ступеней ДЗ от чувствительных ИО при активном сигнале «БК: Срабатывание БК-м».

\item 
Уровень  срабатывания  ИО  по  изменению  тока  обратной  последовательности  задается  уставками $dI2/dt$> и $dI2/dt$\verb|>>|  для чувствительного и грубого канала соответственно. Уровень срабатывания ИО по изменению  тока  прямой  последовательности  задается  уставками  $dI1/dt$> и $dI1/dt$\verb|>>|  для чувствительного и грубого канала соответственно. 

\item
Ввод функции в работу на этапе параметрирования устройства осуществляется программным переключателем SGF1 «Ввод\_функции». Информация о введенной функции отображается выходным сигналом «БК: Ввод».

\item
Функция может быть выведена из работы оперативно путем активации сигнала «ОВ БК», что характеризуется активным состоянием сигнала «БК: Оперативный вывод» на выходе функции. 

\item
Алгоритм работы функции выполнен в соответствии с рисунком \ref{bk:img1}. В таблице \ref{bk:tbl1} приведены параметры, необходимые для настройки функции.

\vspace{1mm}
\begin{figure}[H]
\centering
\includegraphics[width=1\textwidth,height=1\textheight,keepaspectratio]{\fbpath/102. БК/_latex/img/1.pdf}
\captionof{figure}{Функциональная схема логики БК}\label{bk:img1}
\end{figure}


\end{enumerate}